\chapter[Classical sampling and random matrices]{Sampling, spectral concentration, and random-matrix theorems}
\label{ch:classical-sampling}
The following classical theorems connect continuous-time measurement,
spectral estimation, and high-dimensional statistical models. They are
presented here immediately after Fourier and spectral foundations rather
than as a revision-history appendix; each requires its precise hypotheses.

\section{Nyquist--Shannon sampling}
\begin{theorem}[Classical bandlimited reconstruction]
For a finite-energy continuous-time signal with Fourier support
inside $[-B,B]$ Hz, an ideal uniform, noiseless sampling sequence at
$F_s>2B$ determines the signal. With
$\operatorname{sinc}(u)=\sin(\pi u)/(\pi u)$, reconstruction has the form
\[
x(t)=\sum_{n\in\Z}x(n/F_s)\operatorname{sinc}(F_st-n),
\]
in the suitable convergence sense.
\end{theorem}
\emph{Application:} set anti-aliasing and acquisition requirements.
Real filters have finite rolloff; extracellular spikes are not
literally compactly bandlimited.

\section{Bedrosian's theorem}
\begin{theorem}[A sufficient spectral-separation case]
If $a$ has Fourier support in $[-A,A]$ and $c$ is supported
outside $(-A,A)$, with the required transform and product
hypotheses, then $\Hilbert(ac)=a\,\Hilbert(c)$.
\end{theorem}
For $c(t)=\cos(2\pi f_0t)$ with $f_0>A$, an ideal positive envelope
$a(t)$ is recovered from the analytic signal of $a(t)c(t)$.
The nonoverlap condition is sufficient, not a general necessary and
sufficient criterion. PAC remains susceptible to waveform shape,
transients, nonstationarity and filtering \cite{aru2015,cole2019}.

\section{Wiener--Khinchin theorem}
\begin{theorem}
A weakly stationary second-order process has an autocovariance and
a corresponding nonnegative spectral measure. When densities
exist and inversion is justified,
\[
S_x(f)=\int_{\R}R_x(\tau)e^{-2\pi if\tau}\,d\tau,\qquad
R_x(\tau)=\int_{\R}S_x(f)e^{2\pi if\tau}\,df.
\]
\end{theorem}
This supports power spectral estimation; it does not guarantee a unique
oscillatory/aperiodic decomposition \cite{donoghue2020,gerster2022}.

\section{Slepian concentration and DPSS tapers}
For a length-$N$ taper $v$ and normalized half-bandwidth $0<W<1/2$,
define the concentration ratio
\[
\lambda(v)=
\frac{\int_{-W}^W|\sum_{n=0}^{N-1}v_ne^{-2\pi ifn}|^2\,df}
{\int_{-1/2}^{1/2}|\sum_{n=0}^{N-1}v_ne^{-2\pi ifn}|^2\,df}.
\]
\begin{theorem}
Eigenvectors of the discrete spectral-concentration operator
maximize this ratio successively under orthogonality constraints.
Its eigenvalues lie in $[0,1]$.
\end{theorem}
\emph{Application:} \texttt{\detokenize{arsenal/path/multitaper_psd.m}}
uses MATLAB DPSS and averages tapered periodograms, without
adaptive taper weighting.

\section{Johnson--Lindenstrauss lemma}
\begin{theorem}[Finite-set projection]
A suitable normalized random linear map to dimension
$m=O(\varepsilon^{-2}\log P)$ approximately preserves all
pairwise Euclidean distances between $P$ fixed points with high
probability; constants also depend on the desired failure probability.
\end{theorem}
A random dense projection is \emph{not} equivalent to sampling a random
subset of neural coordinates. Coordinate sampling needs additional
coherence, leverage and sampling assumptions. This theorem does not
alone certify reconstruction of neural manifolds.

\section{Marcenko--Pastur and Tracy--Widom}
\begin{theorem}[Marcenko--Pastur law, standard isotropic model]
For large sample covariance matrices from iid isotropic
observations with ratio $N/T\to\gamma$, the continuous
limiting spectrum is supported on
$[\sigma^2(1-\sqrt\gamma)^2,\sigma^2(1+\sqrt\gamma)^2]$;
when $\gamma>1$ there is also an atom at zero.
\end{theorem}
\begin{theorem}[Tracy--Widom edge fluctuation]
For suitable real Wishart ensembles, correctly centered and scaled
largest-eigenvalue fluctuations at the upper edge follow a
Tracy--Widom limit.
\end{theorem}
Neither theorem provides an exact universal noise threshold
for time-correlated, heteroscedastic, or preprocessed neural data.
Edge fluctuations may occur on either side of the limiting edge.

\section{Jensen's inequality}
\begin{theorem}
If $\phi$ is convex and the relevant expectations exist,
$\phi(\E X)\leq\E[\phi(X)]$.
\end{theorem}
An elementary dimensionality bound follows instead directly from
Cauchy--Schwarz:
\[
1\leq D_{\mathrm{PR}}=
\frac{(\sum_{i=1}^r\lambda_i)^2}{\sum_{i=1}^r\lambda_i^2}\leq r,
\quad \lambda_i>0.
\]
Jensen's more direct role in this textbook is the derivation of
variational lower bounds and expectation-maximization, given in the
GLM and optimization chapter. Inferring a structural biological
dimensionality bound requires additional assumptions.

\section{Schoenberg positive-definite radial kernels}
\begin{theorem}[Powered exponential, all Euclidean dimensions]
For $\alpha>0$, the kernel
$k(x,y)=\exp(-\alpha\|x-y\|^q)$
is positive definite on all Euclidean dimensions for
$0<q\leq2$; this all-dimension guarantee does not hold
for $q>2$.
\end{theorem}
This provides a covariance-valid kernel family for Gaussian processes,
but not a proof that a particular model recovers physiological factors.
Matern kernels are a different family with their own assumptions.

\section{Tonelli and Fubini}
\begin{theorem}[Tonelli]
For nonnegative measurable $f$ over a product of sigma-finite measure
spaces, iterated integrals can be interchanged, possibly with value
$+\infty$.
\end{theorem}
\begin{theorem}[Fubini]
For absolutely integrable measurable $f$ on a product of
sigma-finite measure spaces, the integrals can be interchanged
and have finite common value.
\end{theorem}
These theorems justify integrals used in spike-count expectations and
latent-variable models only when their hypotheses apply.
% --- expanded edition (October 2026) ---
\section{The expected periodogram: Wiener--Khinchin in finite samples}
Let $X_t$ be zero-mean stationary with absolutely summable autocovariance
$\gamma$ and spectral density $S(f)=\sum_k\gamma(k)e^{-2\pi ifk}$ (unit
sampling interval). For the periodogram
$I_T(f)=T^{-1}|\sum_{t=0}^{T-1}X_te^{-2\pi ift}|^2$,
\[
\E I_T(f)=\sum_{|k|<T}\Bigl(1-\frac{|k|}T\Bigr)\gamma(k)e^{-2\pi ifk}
=\int_{-1/2}^{1/2}F_T(f-\nu)S(\nu)\,d\nu,
\]
where $F_T(f)=T^{-1}\bigl(\sin(\pi Tf)/\sin(\pi f)\bigr)^2$ is the Fej\'er kernel.
\begin{proof}
Expand the squared modulus into a double sum over $s,t$ and take
expectations: $\E X_s X_t=\gamma(s-t)$. Grouping by $k=s-t$ there are
$T-|k|$ pairs with each lag, giving the first expression. The second
follows by substituting $\gamma(k)=\int S(\nu)e^{2\pi i\nu k}d\nu$ and
recognising $\sum_{|k|<T}(1-|k|/T)e^{2\pi i(\nu-f)k}=F_T(f-\nu)$.
\end{proof}
The periodogram is therefore an unbiased estimate not of $S(f)$ but of
$S$ smoothed by the Fej\'er kernel, whose sidelobes decay only as
$f^{-2}$. Where the spectrum has a large dynamic range, such as the steep
$1/f^\chi$ LFP spectrum, leakage from low frequencies can dominate the
expected periodogram at high frequencies. Tapering replaces $F_T$ by the
squared transform of the taper, which is the rigorous reason that
Slepian tapers reduce bias.

\section{Johnson--Lindenstrauss in numbers}
The lemma guarantees that $n$ points can be mapped linearly to
$k=O(\varepsilon^{-2}\log n)$ dimensions with all pairwise distances
preserved within a factor $1\pm\varepsilon$, and a random Gaussian matrix
scaled by $1/\sqrt k$ achieves this with high probability. A commonly
used explicit sufficient dimension is
$k\ge4\log n/(\varepsilon^2/2-\varepsilon^3/3)$.
\begin{example}[Random projections of population states]
For $n=10^4$ time bins and $\varepsilon=0.2$, the bound gives
$k\ge4\cdot9.21/(0.02-0.00267)\approx2126$. The number does not depend on
the number of recorded neurons. The lemma is therefore a statement that
\emph{distance geometry} of many points survives random compression,
not that neural activity is low-dimensional: the required $k$ grows with
the number of points, whatever the data.
\end{example}

\section{When Fubini fails}
The hypotheses of Fubini's theorem are not decorative. On
$[0,1]^2$, let $f(x,y)=(x^2-y^2)/(x^2+y^2)^2$. Then
$\int_0^1\int_0^1f\,dy\,dx=\pi/4$ while
$\int_0^1\int_0^1f\,dx\,dy=-\pi/4$, because $\int\!\!\int|f|=\infty$. In
spike-train calculations the analogous failure arises when exchanging
an expectation with an infinite sum over spikes whose contributions are
not absolutely summable, for example in shot-noise models with heavy-tailed
kernels.

\section*{Exercises for this chapter}
\begin{exercise}
Show that $F_T(0)=T$ and $\int_{-1/2}^{1/2}F_T(f)\,df=1$.
\end{exercise}
\begin{exercise}
For white noise, show that $\E I_T(f)=\gamma(0)$ exactly for every $T$.
\end{exercise}
\begin{exercise}
Compute the JL dimension for $n=10^6$ and $\varepsilon=0.1$.
\end{exercise}
\begin{exercise}
Verify the two iterated integrals in the Fubini counterexample.
\end{exercise}
